\documentclass[12pt,a4paper]{article}
\usepackage[utf8]{inputenc}
\usepackage[T1]{fontenc}
\usepackage[polish]{babel}
\usepackage{geometry}
\geometry{a4paper, margin=2.5cm}
\usepackage{graphicx}
\usepackage{amsmath}
\usepackage{booktabs}
\usepackage{listings}
\usepackage{xcolor}
\usepackage{hyperref}
\usepackage{enumitem}

\lstset{
    basicstyle=\ttfamily\footnotesize,
    backgroundcolor=\color{gray!10},
    frame=single,
    breaklines=true,
    showstringspaces=false,
    language=Python,
    commentstyle=\color{green!50!black},
    keywordstyle=\color{blue},
    stringstyle=\color{red!70!black},
}

\begin{document}

% ===================== STRONA TYTUŁOWA =====================
\thispagestyle{empty}
\begin{center}
\begin{tabular}{|p{5cm}|p{9cm}|}
\hline
\textbf{Uczelnia} & Politechnika Bydgoska im. J. J. Śniadeckich \\ \hline
\textbf{Wydział} & Telekomunikacji, Informatyki i Elektrotechniki \\ \hline
\textbf{Przedmiot} & Sztuczna Inteligencja \\ \hline
\textbf{Tytuł projektu} & System AI do gry w Kuhn Poker oparty na sieci neuronowej i algorytmie genetycznym \\ \hline
\textbf{Dokument} & Sprawozdanie z postępów -- Tydzień 1 \\ \hline
\textbf{Autorzy} & [Imię i Nazwisko 1] \newline [Imię i Nazwisko 2] \\ \hline
\textbf{Grupa} & 4 \\ \hline
\textbf{Data} & \today \\ \hline
\textbf{Prowadzący} & dr inż. Tomasz Talaśka \\ \hline
\end{tabular}
\end{center}
\vfill
\begin{center}
\Large Rok akademicki 2025/2026
\end{center}
\newpage

% ===================== 1. ZAKRES TYGODNIA =====================
\section{Zakres prac w tygodniu 1}

Celem pierwszego tygodnia było przygotowanie kompletnego szkieletu
projektu: działającego silnika gry, mechanizmu uczenia (algorytm
genetyczny), zestawu przeciwników odniesienia oraz jasno zdefiniowanego
interfejsu pomiędzy częścią Osoby~1 (silnik i ewolucja) a częścią
Osoby~2 (sieć neuronowa). Założeniem było, aby na koniec tygodnia
system dało się uruchomić od początku do końca, nawet przed
ukończeniem właściwej sieci neuronowej.

\subsection*{Podział obowiązków}
\begin{itemize}[nosep]
    \item \textbf{Osoba 1}: silnik gry, kodowanie stanu, system ekspertowy,
          strategia Nasha, algorytm genetyczny, kontrakt agenta, testy
          jednostkowe, dokumentacja (sekcje 1--2).
    \item \textbf{Osoba 2}: implementacja sieci neuronowej zgodnej z
          kontraktem agenta (w trakcie), dokumentacja (sekcje 3--4).
\end{itemize}

% ===================== 2. ZREALIZOWANE ELEMENTY =====================
\section{Zrealizowane elementy (Osoba 1)}

\subsection{Silnik gry}
Zaimplementowano pełne reguły Kuhn Poker w module \texttt{src/engine/kuhn.py}.
Stan gry reprezentowany jest przez kartę gracza (liczba całkowita
1--3) oraz historię akcji (łańcuch znaków, np.\ \texttt{"pb"}).
Zdefiniowano komplet historii terminalnych wraz z wypłatami oraz
funkcję \texttt{play\_hand}, która rozgrywa pojedyncze rozdanie między
dwoma agentami.

\subsection{Kodowanie stanu}
Moduł \texttt{src/engine/state.py} zamienia stan gry na 12-wymiarowy
wektor cech (one-hot karty oraz one-hot trzech pozycji historii).
Dodano funkcje pomocnicze \texttt{build\_state\_dict} oraz
\texttt{encode\_from\_dict} umożliwiające współpracę z enkoderem
Osoby~2 w uzgodnionym formacie słownikowym.

\subsection{Agenci odniesienia}
W module \texttt{src/expert/baseline.py} zaimplementowano:
\begin{itemize}[nosep]
    \item \texttt{random\_agent} -- agent losowy,
    \item \texttt{always\_pass\_agent}, \texttt{always\_bet\_agent} -- agenci skrajni,
    \item \texttt{expert\_agent} -- system ekspertowy oparty na regułach,
    \item \texttt{nash\_agent\_factory} -- strategia równowagi Nasha
          (rodzina parametryzowana wartością $\alpha \in [0, 1/3]$).
\end{itemize}

\subsection{Algorytm genetyczny}
W module \texttt{src/ga/} zaimplementowano funkcję fitness
(\texttt{fitness.py}) oraz pętlę ewolucji opartą na bibliotece DEAP
(\texttt{evolution.py}). Wykorzystano selekcję turniejową, krzyżowanie
typu \emph{blend} oraz mutację gaussowską, z zachowaniem najlepszego
osobnika (Hall of Fame).

\subsection{Kontrakt agenta -- punkt styku z Osobą 2}
Kluczowym elementem tygodnia było ustalenie jednolitego interfejsu
między częściami projektu. W module \texttt{src/agent/base.py}
zdefiniowano kontrakt, który musi spełniać każda klasa agenta:
\begin{enumerate}[nosep]
    \item atrybut klasowy \texttt{WEIGHT\_SIZE} -- długość genomu;
    \item konstruktor \texttt{\_\_init\_\_(self, weights)} -- budowa agenta z wektora wag;
    \item metoda \texttt{\_\_call\_\_(self, card, history) -> Action} -- wywoływalność instancji.
\end{enumerate}
Dzięki wywoływalności instancji agent neuronowy będzie zamiennikiem
\emph{drop-in} dla agentów funkcyjnych, a sama klasa może pełnić rolę
fabryki w algorytmie genetycznym:

\begin{lstlisting}
best, log = run_evolution(NeuralAgent.WEIGHT_SIZE, NeuralAgent, ...)
\end{lstlisting}

Jako referencyjną implementację kontraktu dostarczono klasę
\texttt{LinearAgent} (polityka liniowa), która jednocześnie pełni rolę
tymczasowego zamiennika sieci. Pozwala to uruchomić pełny cykl
treningowy zanim Osoba~2 ukończy sieć neuronową.

\subsection{Testy jednostkowe}
Przygotowano 16 testów (\texttt{pytest}) pokrywających reguły gry,
kodowanie stanu, agentów odniesienia oraz kontrakt agenta. Wszystkie
testy przechodzą.

% ===================== 3. WYNIKI WSTĘPNE =====================
\section{Wyniki wstępne}

Uruchomiono pełny cykl treningu (50 generacji, populacja 60) z
zastępczą polityką liniową. Wytrenowany agent oceniony został w
bezpośrednich pojedynkach (po 2000 rozdań, naprzemienne pozycje):

\begin{center}
\begin{tabular}{lr}
\toprule
Przeciwnik & Średnia wygrana / rozdanie \\ \midrule
agent losowy        & $+0{,}30$ \\
agent zawsze-bet    & $+0{,}40$ \\
system ekspertowy   & $\approx 0$ (wynik zbliżony) \\
strategia Nasha     & $\approx -0{,}06$ \\
\bottomrule
\end{tabular}
\end{center}

Wynik przeciwko strategii Nasha (około $-0{,}06$ żetona na rozdanie)
jest zbieżny z teoretyczną wartością gry $-1/18 \approx -0{,}0556$, co
stanowi wstępne potwierdzenie poprawności silnika oraz funkcji fitness.

% ===================== 4. PLAN NA TYDZIEŃ 2 =====================
\section{Plan na tydzień 2}
\begin{itemize}[nosep]
    \item Osoba 2: implementacja sieci neuronowej (\texttt{NeuralAgent})
          zgodnej z kontraktem; podmiana zastępczej polityki liniowej.
    \item Wspólnie: pierwsze pełne treningi neuroewolucji, strojenie
          hiperparametrów (rozmiar populacji, liczba generacji, $\sigma$ mutacji).
    \item Osoba 1: rozbudowa systemu ekspertowego, przygotowanie
          skryptów do generowania wykresów zbieżności z pliku
          \texttt{logbook.json}.
    \item Dokumentacja: uzupełnienie sekcji 3 (implementacja sieci).
\end{itemize}

% ===================== 5. STAN REPOZYTORIUM =====================
\section{Stan repozytorium}
\begin{lstlisting}[language=bash]
kuhn-poker-ga/
  src/
    engine/   kuhn.py, state.py
    expert/   baseline.py
    agent/    base.py, linear_agent.py, nn_agent_template.py
    ga/       fitness.py, evolution.py
  tests/      test_engine.py, test_agent.py   (16 testow, wszystkie zielone)
  docs/       main.tex, sprawozdanie_tydzien1.tex
  main.py, requirements.txt, README.md
\end{lstlisting}

Uruchomienie: \texttt{pip install -r requirements.txt}, następnie
\texttt{python -m pytest tests/ -v} oraz \texttt{python main.py}.

\end{document}
